\documentclass[10pt,a4paper]{amsart}
\usepackage[margin=19mm]{geometry}
\usepackage{amsmath,amssymb,mathtools}
\usepackage[scr=rsfs,cal=euler]{mathalfa}
\usepackage{xfrac}
\usepackage[unicode,hidelinks,bookmarks=false]{hyperref}
\newcommand{\ie}{i.e.\ }
\newcommand{\cf}{cf.\ }
\newcommand{\resp}{respectively\ }
\newcommand{\Subsection}{\subsection}
\providecommand{\nobreakdash}{\nobreak\hbox{-}}
\setlength{\emergencystretch}{2em}
\multlinegap0pt
\usepackage{amssymb}
\usepackage{enumerate}
\usepackage{tikz}
\newtheorem{thm}{Theorem}
\newtheorem{lem}{Lemma}
\newtheorem{pf}{Pf}
\title{On the logarithmic coefficients of Ma-Minda type convex functions}
\author{Md Firoz Ali, Lokenath Thakur}
\date{Source: arXiv:2603.17940v1 (CC BY 4.0)}
\begin{document}
\maketitle

\noindent\emph{Source and licence.} On the logarithmic coefficients of Ma-Minda type convex functions. Version arXiv:2603.17940v1 (CC BY 4.0), distributed by arXiv under the Creative Commons Attribution 4.0 International licence, http://creativecommons.org/licenses/by/4.0/, as recorded in the arXiv licence metadata for this exact version and shown on the abstract page https://arxiv.org/abs/2603.17940v1. Full licence notice: \texttt{licenses/arxiv-2603.17940-CC-BY-4.0.txt}.\par\smallskip

\noindent\textit{Editorial scope.} Counted unit: the article's thm T-thm-20 (source0.tex lines 461--479). The article's complete original proof of it is delivered with it (source0.tex lines 831--994, one \texttt{proof} environment of 164 lines, which the article places away from the statement). No mathematics is written, repaired or re-derived: the statement and the proof are the article's own lines, in the article's own notation, and no retained line is substituted or re-flowed.  Retained uncounted as the source-local context the proof needs: lines 171--173; lines 571--573; lines 582--612; lines 618--625; lines 822--824.  Nothing was omitted from any retained span, so the package carries no omission record.  The retained lines cite 2 works, whose entries are transcribed verbatim from the article's own bibliography source in the e-print and delivered with short editorial labels recorded in the item's reference mapping.  No retained line contains a figure environment, an image-inclusion command or drawing code, so no image is delivered and the package declares no asset dependency.  Editorial formatting only, in addition: an AMS \texttt{amsart} preamble that restates the article's own declarations and macros used by the retained lines, namely the environments \texttt{thm}, \texttt{lem}, \texttt{pf} on separate counters, together with the standard packages those lines need.  The retained environments print in this document as Theorem 1, Lemma 1, Pf 1 in order of appearance, whereas the article prints the same items with the numbering of its own full document.  The complete native source of the article as distributed in the arXiv e-print for this exact version is bundled unchanged as \texttt{sources/196662/source0.tex}.
\begin{equation}\label{T-05}
f(z) = z + \sum_{n=2}^\infty a_n z^n,
\end{equation}

\begin{align}\label{T-01}
\omega(z)=\sum_{n=1}^\infty c_n z^n.
\end{align}

\begin{lem}\label{T-34}\cite{1981-Prokhorov-szynal}
Let $f\in\mathcal{B}_0$ be of the form \eqref{T-01}. Then for real $\mu$ and $\nu$, we have
\begin{align*}
|c_3+\mu c_1c_2+\nu c_1^3|\le
\begin{cases} 
1 & \text{if } (\mu, \nu) \in D_1 \cup D_2 \cup \{ (2, 1) \} \\
|\nu| & \text{if } (\mu, \nu) \in \bigcup_{k=3}^{7} D_k \\
\frac{2}{3} \left( |\mu| + 1 \right) \left( \frac{|\mu| + 1}{3(|\mu| + 1 +\nu)} \right)^{1/2} & \text{if } (\mu, \nu) \in D_8 \cup D_9 \\
\frac{1}{3} \nu \left( \frac{\mu^2 - 4}{\mu^2 - 4\nu} \right) \left( \frac{\mu^2 - 4}{3(\nu - 1)} \right)^{1/2} & \text{if } (\mu, \nu) \in D_{10} \cup D_{11} - \{ (2, 1) \} \\
\frac{2}{3} \left( |\mu| - 1 \right) \left( \frac{|\mu| - 1}{3(|\mu| - 1 - \nu)} \right)^{1/2} & \text{if } (\mu, \nu) \in D_{12}.
\end{cases}
\end{align*}
where 
\allowdisplaybreaks
\begin{align*}
D_1 &= \left\{ (\mu, \nu): |\mu| \leq \frac{1}{2},\; -1 < \nu \leq 1 \right\}, \\
D_2 &= \left\{ (\mu, \nu): \frac{1}{2} \leq |\mu| \leq 2,\; \frac{4}{27} (|\mu| + 1)^3 - (|\mu| + 1) \leq \nu \leq 1 \right\}, \\
D_3 &= \left\{ (\mu, \nu): |\mu| \leq \frac{1}{2},\; \nu \leq -1 \right\}, \\
D_4 &= \left\{ (\mu, \nu): |\mu| \geq \frac{1}{2},\; \nu \leq -\frac{2}{3} (|\mu| + 1) \right\}, \\
D_5 &= \left\{ (\mu, \nu): |\mu| \leq 2,\; \nu \geq 1 \right\}, \\
D_6 &= \left\{ (\mu, \nu): 2 \leq |\mu| \leq 4,\; \nu \geq \frac{1}{12} (\mu^2 + 8) \right\}, \\
D_7 &= \left\{ (\mu, \nu): |\mu| \geq 4,\; \nu \geq \frac{2}{3} (|\mu| - 1) \right\}, \\
D_8 &= \left\{ (\mu, \nu): \frac{1}{2} \leq |\mu| \leq 2,\; -\frac{2}{3} (|\mu| + 1) \leq \nu \leq \frac{4}{27} (|\mu| + 1)^3 - (|\mu| + 1) \right\}, \\
D_9 &= \left\{ (\mu, \nu): |\mu| \geq 2,\; -\frac{2}{3} (|\mu| + 1) \leq \nu \leq \frac{2|\mu|(|\mu| + 1)}{\mu^2 + 2|\mu| + 4} \right\}, \\
D_{10} &= \left\{ (\mu, \nu): 2 \leq |\mu| \leq 4,\; \frac{2|\mu|(|\mu| + 1)}{\mu^2 + 2|\mu| + 4} \leq \nu \leq \frac{1}{12} (\mu^2 + 2) \right\}, \\
D_{11} &= \left\{ (\mu, \nu): |\mu| \geq 4,\; \frac{2|\mu|(|\mu| + 1)}{\mu^2 + 2|\mu| + 4} \leq \nu \leq \frac{2|\mu|(|\mu| - 1)}{\mu^2 - 2|\mu| + 4} \right\}, \\
D_{12} &= \left\{ (\mu, \nu): |\mu| \geq 4,\; \frac{2|\mu|(|\mu| - 1)}{\mu^2 - 2|\mu| + 4} \leq \nu \leq \frac{2}{3} (|\mu| - 1) \right\}.
\end{align*}

Moreover, all the inequalities are sharp.
\end{lem}

 \begin{lem}\label{T-200}\cite{2016-Iason}
 Let $f\in\mathcal{B}_0$ be of the form \eqref{T-01}. Then for $\lambda \in \mathbb{C} $, we have
\begin{align*}
|c_2+\lambda c_1^2| &\le  \max\{1, |\lambda|\},\\
|c_3+(1+\lambda) c_1c_2+\lambda c_1^3| &\le \max\{1, |\lambda|\}
\end{align*}
and the inequalities are sharp.
 \end{lem}

\begin{align}\label{T-15}
\gamma_1=\frac{1}{2}a_2,\quad \gamma_2=\frac{1}{2}\left(a_3-\frac{1}{2}a_2^2\right),\quad \gamma_3=\frac{1}{2}\left(a_4-a_2a_3+\frac{1}{3}a_2^3\right).
\end{align}

\begin{thm}\label{T-thm-20}
Let \( f \in \mathcal{C}(A,B) \) with \( -1 \leq B < A \leq 1 \). Then the logarithmic coefficients \( \gamma_n \) of \( f(z) \) satisfy
\begin{align*}
|\gamma_1| &\leq \frac{A - B}{4},\\
|\gamma_2| &\leq 
\begin{cases}
\displaystyle \frac{A-B}{12} & ~~\text{for}~ |A-5B|\le 4,\\[2mm]
\displaystyle \frac{A-B}{48}|A-5B| & ~~\text{for}~ |A-5B|>4,
 \end{cases}\\[3mm]
|\gamma_3| &\leq 
\begin{cases}
\displaystyle \frac{A - B}{24}, &~~\text{for}~(\mu,\nu)\in D_1\cup D_2, \\[2mm]
\displaystyle \frac{A - B}{48}|B(A - 3B)|, & ~~\text{for}~(\mu,\nu)\in D_6, \\[2mm]
\displaystyle \frac{A-B}{72\sqrt{3}} \cdot \frac{\left( |A - 5B| + 2 \right)^{3/2}}{\sqrt{|A - 5B| + 2 + 3B^{2} - AB}}  & ~~\text{for}~(\mu,\nu)\in D_8\cup D_9,
\end{cases}
\end{align*}
where \(\mu = \frac{1}{2}(A - 5B)\) and \(\nu = \frac{1}{2}(3B^2 - AB)\) and the sets $D_1,D_2, D_6, D_8,D_9$ are defined as in Lemma \ref{T-34}.
All the estimates are sharp.
\end{thm}

\begin{pf}[\textbf{Proof of Theorem \ref{T-thm-20}}]
Let $ f \in \mathcal{C}(A,B) $ be of the form \eqref{T-05} . Then
$$
1 + \frac{zf''(z)}{f'(z)} \prec \frac{1+Az}{1+Bz}.
$$  
Thus, there exist an $\omega\in\mathcal{B}_0$ of the form \eqref{T-01} such that
$$
1 + \frac{zf''(z)}{f'(z)} = \frac{1+A \omega(z)}{1+B \omega(z)}.
$$
Now equating the coefficients of $z^2,z^3$ and $z^4$, and simplifying, we obtain
    \begin{align*}
    a_2&=\frac{c_1(A-B)}{2}, \quad a_3=\frac{A-B}{6}\left((A-2B)c_1^2+c_2\right),\\
    a_4&= \frac{A-B}{24}\left(2c_3+(3A-7B)c_{1}c_{2}+(A^2-5AB+6B^2)c_1^3\right).
    \end{align*}
Substituting these values into \eqref{T-15}, we obtain
\begin{align}
    \gamma_1&=\frac{c_1(A-B)}{4}, \quad \gamma_2=\frac{A-B}{48}\left((A-5B)c_1^2+4c_2\right),\label{T-91}\\
    \gamma_3&= \frac{A-B}{24}\left(c_3+\frac{1}{2}(A-5B)c_{1}c_{2}+\frac{1}{2}(3B^2-AB)c_1^3\right).\label{T-92}
    \end{align}
Hence from \eqref{T-91}, we find
    $$
     |\gamma_1| \leq \frac{A - B}{4}
    $$
    and the inequality is sharp for the function $g_1\in\mathcal{C}(A,B)$ defined by
\begin{align}\label{T-98}
1 + \frac{zg_1''(z)}{g_1'(z)} = \frac{1+Az}{1+B z}.
\end{align}
Furthermore, from \eqref{T-91} and Lemma \ref{T-200}, we get
\begin{align*}
|\gamma_2|
&=\frac{A-B}{12}\left|c_2+\frac{A-5B}{4}c_1^2\right|\\
&\le \frac{A-B}{12}\max\left\{1,\frac{|A-5B|}{4}\right\}\\
& =\begin{cases}\frac{A-B}{12} & ~~\text{for}~ |A-5B|\le 4\\[2mm]
\frac{A-B}{48}|A-5B| & ~~\text{for}~ |A-5B|>4
 \end{cases}
\end{align*}
The first inequality is sharp for the function $g_2\in\mathcal{C}(A,B)$ defined by
\begin{align}\label{T-93}
1 + \frac{zg_2''(z)}{g_2'(z)} = \frac{1+Az^2}{1+Bz^2}.
\end{align}
For the function $g_2\in\mathcal{C}(A,B)$, it is easy to see that
$$
a_2=0,\quad a_3=\frac{A-B}{6}\quad\text{and}\quad\gamma_2(g_2)=\frac{1}{2}\left(a_3-\frac{1}{2}a_2^2\right)=\frac{A-B}{12}.
$$
The second inequality is sharp for the function $g_1\in\mathcal{C}(A,B)$ defined by \eqref{T-98}. For the function $g_1\in\mathcal{C}(A,B)$, it is easy to see that
$$
a_2=\frac{A-B}{2},\quad a_3=\frac{A^2-3AB+2B^2}{6}
$$
and so 
$$\gamma_2(g_1)=\frac{1}{2}\left(a_3-\frac{1}{2}a_2^2\right)=\frac{A-B}{48}(A-5B).$$




Further, from \eqref{T-92}, we get
\begin{align}\label{T-905}
|\gamma_{3}|
&=\frac{A-B}{24}\left|c_3+\frac{1}{2}(A-5B)c_{1}c_{2}+\frac{1}{2}(3B^2-AB)c_1^3\right|\\
&= \frac{A-B}{24}\left|c_3+\mu c_{1}c_{2}+\nu c_1^3\right|\nonumber,
\end{align}
with \(\mu = \frac{1}{2}(A - 5B)\) and \(\nu = \frac{1}{2}(3B^2 - AB)\). We now use Lemma \ref{T-34} to find the estimate of $\gamma_3$. For \(-1 \leq B < A \leq 1\), one can quickly verify that 
\begin{equation}\label{T-900}
\mu \in (-2, 3]\quad \text{and}\quad \nu \in \left[-\frac{1}{24}, 2\right].
\end{equation}

Suppose that the sets $D_1,D_2,\ldots,D_{12}$ are defined as in Lemma \ref{T-34}.
\begin{itemize}
\item Since \(|\mu| \le 4\), it immediately follows that the sets \(D_7\), \(D_{11}\), and \(D_{12}\) are empty. Similarly, $D_3$ is empty because $\nu \not\le -1$.

\item If $|\mu|\ge \frac{1}{2}$ then
$$-\frac{2}{3}(|\mu| + 1)\le -1 $$
and so the condition $\nu \leq -\frac{2}{3}(|\mu| + 1)$ is incompatible with the admissible range of \(\nu\) given in \eqref{T-900}. Hence, the set $D_4$ is empty.

\item When $2\le |\mu|\le 4$ we must have $\mu\in(2,3]$. In this case,
$$
\frac{2|\mu|(|\mu| + 1)}{\mu^2 + 2|\mu| + 4}> \frac{2(\mu^2 + \mu)}{9 + 6 + 4}> \frac{\mu^2 + 2}{12}.
$$
This shows that \(D_{10}\) is also empty.

\item The conditions for $D_5$ ($|\mu| \le 2$ and $\nu \ge 1$) is equivalent to 
\[
\frac{A-4}{5} \leq B \leq \frac{A - \sqrt{A^2 + 24}}{6}, \quad |A| < 1.
\]
However, for any \(|A| < 1\), we have \(\frac{A-4}{5} > \frac{A - \sqrt{A^2 + 24}}{6}\). Hence, the set \(D_5\) is also empty.

\end{itemize}

On the other hand, it is easy to check that 
\begin{itemize}
\item when \((A,B)=(0,-0.1)\) then the associated $(\mu,\nu)\in D_1$,
\item when \((A,B)=(-0.5,-0.6)\) then the associated $(\mu,\nu)\in D_2$,
\item when \((A,B)=(0,-0.9)\) then the associated $(\mu,\nu)\in D_6$,
\item when \((A,B)=(0,-0.795)\) then the associated $(\mu,\nu)\in D_8$,
\item when \((A,B)=(0,-0.81)\) then the associated $(\mu,\nu)\in D_9$.
\end{itemize}
This shows that the sets $D_1, D_2, D_6, D_8$ and $D_9$ are non empty and together they cover all possible values of $(A,B)$ satisfying \(-1 \leq B < A \leq 1\).
Now we consider the following cases:\\


\textbf{Case-1:}  Let $(\mu,\nu)\in D_1\cup D_2$. Then from \eqref{T-905} and by Lemma \ref{T-34}, we have
$$
|\gamma_{3}|\le\frac{A-B}{24}.
$$
The inequality is sharp for the function $g_3\in\mathcal{C}(A,B)$ defined by
$$
1 + \frac{zg_3''(z)}{g_3'(z)} = \frac{1+Az^3}{1+Bz^3}.\\
$$
For the function $g_3\in\mathcal{C}(A,B)$, it is easy to see that $a_2=0,~~ a_3=0,~~ a_4=\frac{A-B}{12}$ and so
$$
\gamma_3(g_3)=\frac{1}{2}\left(a_4-a_2a_3+\frac{1}{3}a_2^3\right)=\frac{A-B}{24}.
$$



\textbf{Case-2:} Let $(\mu,\nu)\in D_6$. Then from \eqref{T-905} and by Lemma \ref{T-34}, we have
$$
|\gamma_{3}|\le\frac{A-B}{48}|B(A-3B)|.
$$
The inequality is sharp for the function $g_1\in\mathcal{C}(A,B)$ defined by defined by \eqref{T-98}. For the function $g_1\in\mathcal{C}(A,B)$, it is easy to see that
$$
a_2=\frac{A-B}{2},~~ a_3=\frac{A^2-3AB+2B^2}{6}, ~~a_4=\frac{A^3-6A^2B+11AB^2-6B^3}{24}
$$
and so 
$$\gamma_3(g_1)=\frac{1}{2}\left(a_4-a_2a_3+\frac{1}{3}a_2^3\right)=-\frac{A-B}{48}B(A-3B).$$



\textbf{Case-3:} Let $(\mu,\nu)\in D_8\cup D_9$. Then from \eqref{T-905} and by Lemma \ref{T-34}, we have
$$
|\gamma_3|
\le \frac{A-B}{24}\cdot \frac{2\left( |\mu| + 1 \right)}{3}  \left( \frac{|\mu| + 1}{3(|\mu| + 1 +\nu)} \right)^{1/2}
= \frac{A-B}{72\sqrt{3}} \cdot \frac{\left( |A - 5B| + 2 \right)^{3/2}}{\sqrt{|A - 5B| + 2 + 3B^{2} - AB}}.
$$
The inequality is sharp for the function $g_4\in\mathcal{C}(A,B)$ defined by
$$
1 + \frac{zg_4''(z)}{g_4'(z)} = \frac{1+A\omega(z)}{1+B\omega(z)}
$$
where
\begin{align*}
\omega(z)
&=\frac{z(c-\operatorname{sgn} (\mu)~ z)}{(1-\operatorname{sgn} (\mu)~cz)}, \quad c=\left( \frac{|\mu|+1}{3(|\mu|+\nu+1)}\right)^{1/2}\\[3mm]
&= cz+(c^2-1)\operatorname{sgn} (\mu) z^2+c(c^2-1)z^3+\cdots.
\end{align*}
Here we note that if $(\mu,\nu)\in D_8\cup D_9$ then
$$
3(|\mu|+\nu+1)\ge 3(|\mu|+1)-\frac{2}{3}(|\mu|+1)> |\mu|+1>0
$$
and consequently $|c|<1$. For the function $g_4\in\mathcal{C}(A,B)$, from \eqref{T-905}, we have
%\begin{align*}
%a_2 &= \frac{c(A-B)}{2}, ~~a_3 = \frac{A - B}{6} \left( (A - 2B + 1) c^2 - 1 \right), \\[4pt]
%a_4 &= \frac{A - B}{24} \, c \Bigl( (-3A + 7B - 2) + \bigl( A^2 - 5AB + 3A + 6B^2 - 7B + 2 \bigr) c^2 \Bigr),
%\end{align*}
\begin{align*}
|\gamma_3(g_4)|
&= \frac{A-B}{24}\left|c_3+\mu c_{1}c_{2}+\nu c_1^3\right|\\
&=\frac{A-B}{24}\left|c(c^2-1)+\mu c(c^2-1)\operatorname{sgn} (\mu) +\nu c^3\right|\\
&=\frac{c(A-B)}{24}\left|c^2(1+|\mu|+\nu)-(|\mu|+1)\right|\\
&=\frac{c(A-B)}{24}\left|\frac{1}{3}(|\mu|+1)-(|\mu|+1)\right|\\
&=\frac{c(A-B)}{36}(|\mu|+1)\\
&= \frac{A-B}{72\sqrt{3}} \cdot \frac{\left( |A - 5B| + 2 \right)^{3/2}}{\sqrt{|A - 5B| + 2 + 3B^{2} - AB}}.
\end{align*}


\end{pf}

\begin{thebibliography}{99}
\bibitem[1981-P]{1981-Prokhorov-szynal} {\sc D. V. Prokhorov and J. Szynal}, Inverse coefficients for $(\alpha, \beta)$-convex functions, \textit{Ann. Univ. Mariae Curie-Sk\l odowska, Sect. A}, \textbf{35}~(1981), 125--143.
\bibitem[2016-I]{2016-Iason} {\sc I. Efraimidis}, A generalization of Livingston’s coefficient inequalities for functions with positive real part, \textit{J. Math. Anal. Appl.}, \textbf{435} (2016), 369--379.
\end{thebibliography}
\end{document}
